\section*{Resumen}

El trabajo presenta el diseño, la implementación y la validación de un analizador léxico para un subconjunto de Prolog. La solución, desarrollada en Python, recorre la fuente de izquierda a derecha y reconoce átomos, variables, números, cadenas, operadores y delimitadores. Los espacios y comentarios se descartan sin perder el seguimiento de línea y columna. Para resolver lexemas con prefijos compartidos se aplica máxima coincidencia; además, los operadores alfabéticos \texttt{is} y \texttt{mod} se clasifican después de consumir el identificador completo. Cada token conserva su tipo, lexema, posición y, cuando corresponde, un índice en la tabla de lexemas. El signo de los números se trata como operador independiente, de modo que no se confunde el reconocimiento léxico con la interpretación sintáctica. La validación incluye 40 pruebas automatizadas, un programa correcto y otro con errores recuperables. Las pruebas finalizaron sin fallos; el archivo válido produjo 94 tokens y ningún error, mientras que el archivo inválido produjo 30 tokens y seis diagnósticos. El alcance no incluye análisis sintáctico, semántica ni compatibilidad completa con ISO Prolog.

\textit{Palabras clave: análisis léxico, Prolog, autómata finito, máxima coincidencia, tabla de lexemas.}

\section{Introducción}

El análisis léxico transforma una secuencia de caracteres en unidades que pueden ser tratadas por las etapas posteriores de un compilador o intérprete. En este trabajo, el problema consiste en recorrer archivos escritos en un subconjunto de Prolog, distinguir las categorías definidas y asociar a cada token su posición inicial. El recorrido también debe descartar comentarios y espacios sin alterar el cálculo de línea y columna, además de informar errores léxicos sin detenerse ante el primer carácter inválido.

La dificultad no se reduce a comparar caracteres con una lista. Varias categorías comparten prefijos: \verb|\==| comienza como \verb|\=|, \verb|=..| comienza como \verb|=| y el átomo \texttt{island} contiene el prefijo \texttt{is}. Si el analizador aceptara el primer patrón disponible, fragmentaría lexemas que pertenecen a una sola unidad. Por esa razón, la solución conserva la aceptación provisional más larga y aplica una prioridad estable cuando dos reglas pueden reconocer la misma entrada, criterio habitual en la construcción de analizadores léxicos \cite{dragon}.

El informe define primero el subconjunto admitido y sus expresiones regulares. Luego relaciona esas definiciones con AFD para identificadores, números y operadores, e incluye una determinización y una minimización representativas. Las secciones posteriores describen la arquitectura, la tabla de lexemas, el tratamiento de errores y los resultados del corpus automatizado. De este modo, la especificación formal y el comportamiento del programa pueden contrastarse sobre los mismos casos.
